\bookchapter{Глава 35: Дети и смартфоны}

\chapterintro{В этой главе мы не станем ни рисовать смартфон «цифровым злом», ни превозносить принцип «чем раньше начнёшь — тем умнее». Мы разберём всё по возрастным этапам: когда выдавать устройство, какие вводить правила, как прививать цифровую грамотность и беречь здоровье — шаг за шагом.}

\booksection{В каком возрасте давать ребёнку телефон?}

Универсального ответа нет, но есть возрастные ориентиры. Если дать слишком рано — у ребёнка ещё не дозрела префронтальная кора, самоконтроль не поспевает за «зависимым дизайном» алгоритмов; слишком поздно — и он рискует оказаться в «социальной изоляции» среди сверстников. Ниже — рекомендации по возрастам:

\begin{itemize}
  \item \textbf{3–6 лет}: по возможности не давать смартфон. Если пришлось — не более 15 минут за раз, только образовательный контент, отобранный родителями, и обязательно\textbf{смотреть вместе}. Именно на этом этапе экран сильнее всего влияет на зрение и развитие внимания.
  \item \textbf{7–9 лет}: можно «брать в долг» у родительского устройства — для онлайн-уроков, аудиосказок, поиска информации. Формируется правило «попросил перед использованием — вернул после»; собственный телефон не нужен.
  \item \textbf{10–12 лет}: можно рассмотреть «телефон без интернета» или детские часы — для связи по дороге в школу. Если действительно нужен смартфон, включайте строгий родительский контроль.
  \item \textbf{13 лет и старше}: переход в среднюю школу, растут социальные потребности — смартфон уместен, но вместе с «пользовательским договором» и регулярным разбором использования.
\end{itemize}

\begin{tcbinfo}{Принцип «достаточно — значит достаточно»}
Первый телефон ребёнку не обязан быть флагманом. «Избыточная производительность» процессора чаще всего уходит на игры. Лучше выбрать аппарат с ёмкой батареей, средними характеристиками и экраном, бережным к глазам — в диапазоне 1500–2500 юаней или официальный восстановленный. Им проще управлять. Ориентир по бюджету: 1500–2500 юаней покрывают потребности большинства школьников.

\end{tcbinfo}

\booksubsection{Умные часы как переходное решение}

Если нужно просто «дозвониться после школы», детские часы подойдут лучше телефона. Они только звонят, отправляют голосовые сообщения и показывают геолокацию — без браузера и коротких видео, это естественно отсекает «точки входа в зависимость»:

\begin{itemize}
  \item \textbf{Часы Xiaotiancai}: самая полная экосистема, «добавить друга касанием» — ключевая фишка для общения сверстников, но экосистема приложений закрытая.
  \item \textbf{Детские часы Huawei}: хорошая связь с телефонами Huawei, точное позиционирование, поддержка видеозвонков, выгодное соотношение цена/качество.
  \item \textbf{Детские часы Xiaomi}: доступная цена, полный базовый набор — для семей с ограниченным бюджетом.
\end{itemize}

При покупке часов обратите внимание: выбирайте модель с поддержкой\textbf{4G всех операторов}(2G/3G постепенно выводятся из эксплуатации), и уточните тариф детского пакета у оператора.

\booksection{Полный гид по родительскому контролю}

Прежде чем отдать устройство ребёнку, первое действие — не включение, а настройка родительского контроля. У обеих платформ есть собственные решения:

\booksubsection{iOS: Экранное время}

\begin{itemize}
  \item \textbf{Лимиты на приложения}: Настройки > Экранное время > Лимиты на приложения — для каждой категории приложений можно задать суточное время (например, 30 минут на игры).
  \item \textbf{Время отключения}: задаёт ежедневный «режим тишины» (например, 22:00–7:00); в этот период все приложения становятся серыми и недоступны, доступен только телефон.
  \item \textbf{Ограничения контента}: фильтрация фильмов, веб-сайтов и установки приложений по возрасту, а также запрет покупок внутри приложений.
  \item \textbf{Семейный доступ}: удалённое управление устройством ребёнка с родительского iPhone — одобрение запросов на загрузку и просмотр отчётов об использовании.
\end{itemize}

\booksubsection{Android: Google Family Link / фирменный родительский контроль}

\begin{itemize}
  \item \textbf{Google Family Link}: официальное решение Google — удалённое управление временем использования, блокировка устройства, просмотр отчётов об активности. На китайских Android-смартфонах доступность зависит от модели.
  \item \textbf{Детское пространство Xiaomi}: Настройки > Детское пространство — задаёт время использования и белый список приложений, выход по родительскому паролю.
  \item \textbf{Студенческий режим Huawei}: Настройки > Здоровое использование телефона > Студенческий режим — ограничение временных интервалов и приложений, поддержка удалённого управления.
  \item \textbf{Детский режим OPPO/vivo}: включается в «Настройках» или в фирменном приложении «Детское пространство»; принцип аналогичен.
\end{itemize}

\begin{center}
\begin{tabularx}{\linewidth}{|X|X|X|}
\hline
\textbf{Функция} & \textbf{iOS Экранное время} & \textbf{Android Family Link / фирменные решения} \tabularnewline
\hline
Лимиты на приложения & Да, по категориям & Да, по приложениям \tabularnewline
Время отключения & Да & Да (между брендами есть отличия) \tabularnewline
Фильтрация контента & Системная, более полная & Зависит от бренда, неоднородно \tabularnewline
Удалённое управление & Семейный доступ & Family Link / облако бренда \tabularnewline
Согласование установки приложений & Требуется одобрение родителя & Поддерживается в Family Link \tabularnewline
\hline
\end{tabularx}
\end{center}

\begin{tcbwarn}{Родительский контроль — не «сделал и забыл»}
Дети всегда найдут способ обойти ограничения (в сети полно инструкций «как взломать Экранное время»). Технические средства — лишь базовая линия, настоящая защита — это общение и доверие между родителем и ребёнком. Рекомендуется ежемесячно вместе с ребёнком разбирать отчёты об использовании, а не следить односторонне.

\end{tcbwarn}

\booksection{Режим для подростков}

У всех основных контентных платформ встроен «режим для подростков»: после включения ограничивается время использования, фильтруется нежелательный контент, отключаются платежи и чаевые. Это обязательная функция по закону, но\textbf{родитель должен сам её включить и задать пароль}:

\begin{itemize}
  \item \textbf{Режим для подростков Douyin/Kuaishou}: не более 40 минут в день (запрет 22:00–6:00), показывается только отобранный научно-популярный и образовательный контент; прямые эфиры, чаевые и платежи невозможны.
  \item \textbf{Режим для подростков WeChat}: ограничение WeChat Pay (требуется родительская квота), ограничение официальных аккаунтов и видеороликов; функции «встряхнуть» и «поблизости» по умолчанию отключены.
  \item \textbf{Система противозависимости в играх}: национальная аутентификация по реальному имени + ограничения по времени. Несовершеннолетние могут входить в игры только в пятницу, субботу, воскресенье и праздничные дни с 20:00 до 21:00, по 1 часу в день.
\end{itemize}

\textbf{Не используйте удостоверение родителя для «переписывания» аккаунта}\par

Многие дети регистрируют игровой аккаунт на удостоверение родителя, чтобы обойти противозависимую систему. Это нарушает пользовательское соглашение и лишает ребёнка защиты. Родители должны беречь свои платёжные пароли и личные данные и не «одалживать аккаунт» ради удобства.

\booksection{Цифровая грамотность}

Технологии отсекают 90\% рисков, остальные 10\% — на усмотрение ребёнка. Важнее не запретить интернет, а научить им пользоваться:

\begin{itemize}
  \item \textbf{Защита приватности}: не сообщать незнакомцам настоящее имя, школу, домашний адрес; не отправлять личные фото, удостоверение, семейные снимки; геолокация по умолчанию выключена, включается только при необходимости.
  \item \textbf{Распознавание и реакция на кибербуллинг}: оскорбления, бойкот, вредоносные фотожабы — всё это буллинг. Научите ребёнка не отвечать оскорблениями, не удалять доказательства, сразу сообщить родителям. В WeChat можно долгим нажатием подать жалобу, платформа накажет нарушителя.
  \item \textbf{Критика контента}: учим различать «факт» и «мнение», распознавать клибейт и шокирующие заголовки; понимать, что алгоритм подсказывает лишь то, что вы хотите видеть, а не всю картину мира; не верить рекламе «бесплатный сыр».
  \item \textbf{Кибербезопасность}: не переходить по незнакомым ссылкам, не сканировать QR-коды неизвестного происхождения, не скачивать приложения не из официальных магазинов; пароль не сообщать никому (даже «лучшему другу»); любой, кто просит код подтверждения, — мошенник.
\end{itemize}

\begin{tcbpractice}{Практика: вместе с ребёнком подпишите «Договор об использовании смартфона»}
Записанные правила, подписанные обеими сторонами и прикреплённые на видном месте, работают лучше устных договорённостей. Рекомендуется включить:

\begin{itemize}
  \item Временные интервалы: например, «доступно с 19:30 до 20:30 после выполнения домашних заданий».
  \item Зоны запрета: за обеденным столом, в спальне перед сном, на уроках телефон убирается.
  \item Границы контента: не устанавливать приложения, о которых не знают родители; не вести переписку с незнакомцами.
  \item Последствия нарушения: например, «одно превышение — на следующий день время использования сокращается вдвое».
  \item Обязательства родителя: не «залипать» в телефоне при ребёнке дольше оговорённого времени, подавать личный пример.
\end{itemize}

\end{tcbpractice}

\booksection{Привычки здорового использования}

Глаза, шейный отдел позвоночника и сон ребёнка находятся в стадии развития; телефон вредит здоровью сильнее, чем взрослым:

\begin{itemize}
  \item \textbf{Рекомендации по экранному времени}: 3–6 лет — не более 1 часа в день; 6–12 лет — образовательный контент не ограничивается, развлекательный — не более 1 часа; старше 12 лет — по договорённости с ребёнком, но развлекательный — не более 2 часов.
  \item \textbf{Гигиена зрения}: правило «20-20-20» — каждые 20 минут смотрения экрана смотреть 20 секунд на объект в 20 футах (около 6 метров) вдаль; за час до сна не смотреть в телефон, синий свет подавляет мелатонин и мешает сну.
  \item \textbf{Контроль позы}: экран на уровне глаз, не опускать голову и не сутулиться; после 30 минут непрерывного использования — встать и размяться.
  \item \textbf{Общие семейные правила}: в столовой, спальне, кабинете — «зоны без телефонов»; при семейном просмотре кино и активном отдыхе на улице телефоны убираются; «отложить телефон» — часть семейной культуры, а не наказание для ребёнка.
\end{itemize}

\begin{tcbinfo}{Родитель — лучшая «брандмауэр»}
Исследования показывают сильную положительную корреляцию между временем использования телефона у родителей и у детей. Требовать от ребёнка не смотреть в телефон, а самому листать короткие видео до поздней ночи — правила не сработают. Чтобы изменить ребёнка, начните с того, чтобы отложить собственный телефон.

\end{tcbinfo}

\booksection{Цифровое наследие и дети}

Эту тему легко упустить, но она важна. Цифровые активы ребёнка (аккаунты, фото, остатки на счетах, игровое снаряжение) также сталкиваются с вопросами наследования и приватности:

\begin{itemize}
  \item \textbf{Наследование цифровых активов ребёнка}: в случае трагической смерти несовершеннолетнего процесс наследования социальных аккаунтов, игровых аккаунтов и облачных фото такой же, как у взрослых — требуется свидетельство о смерти и подтверждение родства. Родителям рекомендуется заранее настроить для важных аккаунтов ребёнка «контактное лицо по наследству» или «экстренный контакт» (например, контакт наследника Apple, передача аккаунта WeChat).
  \item \textbf{Цифровая приватность несовершеннолетних}: закон «О защите персональных данных» требует согласия опекуна для обработки персональных данных детей до 14 лет и разработки отдельных правил приватности. Родителям следует регулярно проверять политику приватности приложений, часто используемых ребёнком, и отключать переключатели «персонализированные рекомендации» и «обмен данными».
  \item \textbf{Рекомендации по регистрации аккаунтов}: аккаунты для ребёнка следует регистрировать на контролируемый родителями адрес электронной почты/номер телефона — для единого управления на этапе несовершеннолетия; после совершеннолетия аккаунт можно передать.
\end{itemize}

\begin{tcbwarn}{Остерегайтесь обратной стороны «цифрового наследия детей»}
Ребёнка «записывают» с рождения — фотографии роста, выкладываемые родителями, видео из детского сада, данные о здоровье с умных часов. Эти данные после загрузки сложно полностью удалить. Перед публикацией подумайте: не будет ли ребёнку неловко за это фото через десять лет? Уважайте цифровую личность ребёнка — начните с умеренности в публикациях о детях.

\end{tcbwarn}

\begin{tcbsummary}{Итоги главы}
\begin{itemize}
  \item Нет универсального возраста для выдачи телефона, но в 3–6 лет лучше не давать, а до 13 лет — переход через детские часы или телефон без интернета.
  \item «Достаточно — значит достаточно»: первый телефон — это аппарат за 1000 юаней, а не флагман; избыточная производительность равна избыточным соблазнам.
  \item Экранное время iOS и Family Link/фирменный контроль Android — две технические линии защиты, надёжнее использовать вместе.
  \item Режим для подростков и противозависимая система в играх обязательны по закону, но родитель должен сам включить их и беречь пароль.
  \item Цифровая грамотность важнее технической блокировки: учите ребёнка распознавать буллинг, защищать приватность и критически оценивать контент.
  \item Здоровые привычки начинаются с родителя — вы откладываете телефон, и ребёнок отложит свой.
  \item Цифровое наследие и приватность детей — новая тема; заранее настройте «контакт по наследству» и будьте сдержанны в публикациях о детях.
\end{itemize}

\end{tcbsummary}
